\section*{Public-web capture controls}

The public-web path is built and locally tested. It opens only after a public case has a reviewed
Companies House number and a fixed reporting cut-off. The Responses API then runs
\texttt{web\_search} with \texttt{store=False} to plan discovery, and the application admits only
canonical HTTPS addresses that also appear in the tool's own list of consulted sources. Page titles
and model prose stay metadata; they never become company facts. Application code, not the search
result, then fetches each candidate. It refuses embedded credentials, non-default ports, local hosts
and addresses that do not resolve publicly; it checks robots rules, re-validates every redirect,
permits six textual content types, and enforces limits on elapsed time, redirects and bytes. What
survives capture is converted to visible text, stripped of personal contact details, and stored
under a content hash in a directory that refuses to overwrite. A proposed statement is admitted only
if its address matches a captured source, the statement equals the evidence passage, and that
passage occurs literally in the stored text. Four tasks run in order: discover, capture, extract
and compose, under one fingerprint covering identity, cut-off, model, prompt, policy and budgets.
Table~\ref{tab:sys-failure-recovery-status} separates tested engineering work from work that is on
hold. Retryable tasks keep their attempts and final reasons, and interrupted work resumes only
under the same run contract.

This appendix describes the implemented controls. Chapter~5 reports only the controlled tests of
them. No authorised live company or open website was contacted.
